\section*{Overview}

BFS and DFS are the first graph routines worth internalizing. They are simple enough to write from memory and strong
enough to sit inside many larger algorithms later: topological sort, SCC, bridges, tree DP, shortest paths on
unweighted graphs, and more.

\section*{When to Use It}

Use BFS or DFS when:

\begin{itemize}
  \item the graph structure itself is the problem,
  \item you need reachability, connected components, or a traversal order,
  \item the graph is a tree and you need parent/depth information.
\end{itemize}

\section*{Core Idea}

DFS follows one path as far as it can before backtracking. BFS expands in layers from the starting node. Both visit
every reachable vertex, but the order they expose is different and useful in different ways.

\section*{Key Insight}

BFS is the right tool when distance in number of edges matters. DFS is the right tool when the recursion tree or entry /
exit structure matters. A lot of graph problems become easy once that distinction is clear.

\section*{Operations / Main Technique}

\begin{itemize}
  \item BFS from one source for unweighted shortest paths,
  \item DFS for components, subtree structure, or cycle detection,
  \item repeated traversal from unvisited nodes for a full graph decomposition.
\end{itemize}

\paragraph{Worked example.}
On an unweighted graph, BFS from node \(s\) reaches every vertex at distance \(1\) before any vertex at distance \(2\).
That is exactly why the first time a node is popped, its distance is final.

\section*{Correctness Intuition}

DFS visits every edge leaving a reached vertex, so it eventually explores the whole reachable region. BFS keeps a queue
ordered by layer, so nodes are expanded in nondecreasing distance from the source.

\section*{Complexity Analysis}

With adjacency lists, both traversals run in \(O(n + m)\) time and \(O(n)\) extra memory.

\section*{Implementation}

The code includes:

\begin{itemize}
  \item BFS distances on an adjacency list,
  \item iterative DFS order collection,
  \item a connected-components helper.
\end{itemize}

\section*{Common Pitfalls}

\begin{itemize}
  \item Forgetting to clear the visited array between test cases.
  \item Recursive DFS stack overflow on very deep graphs.
  \item Using BFS when weighted edges invalidate the layer guarantee.
  \item Mixing 0-indexed and 1-indexed graph input.
\end{itemize}

\section*{Variants / Extensions}

\begin{itemize}
  \item Multi-source BFS.
  \item 0-1 BFS for edge weights \(0\) and \(1\).
  \item DFS timestamps for bridges, SCC, or subtree intervals.
\end{itemize}

\section*{Practice Problems}

\begin{itemize}
  \item Count connected components.
  \item Shortest path in an unweighted graph.
  \item Check whether a graph is bipartite with BFS coloring.
  \item Gather subtree sizes in a rooted tree with DFS.
\end{itemize}

\section*{References}

\begin{itemize}
  \item \href{https://cp-algorithms.com/graph/breadth-first-search.html}{cp-algorithms: Breadth First Search}.
  \item \href{https://cp-algorithms.com/graph/depth-first-search.html}{cp-algorithms: Depth First Search}.
  \item \href{https://usaco.guide/silver/graph-traversal}{USACO Guide: Graph Traversal}.
  \item \href{https://codeforces.com/catalog}{Codeforces Catalog}.
\end{itemize}
